\begin{description}
\item[(a)] From the Wien's displacement law the temperature of a black body is
  related to a peak wavelength of emission by the relation:
  \[ \lambda_{max} T = b \]
  \hfill(1.0)
  \begin{align*}
    T_A &= 5978\,\mathrm{K} \\
    T_B &= 4830\,\mathrm{K}
  \end{align*}
  \hfill(0.5 + 0.5)

\item[(b)] Due to the diffraction minimal angular separation between two
  celestial objects to distinguish them:
  \[ \delta = 1.22\,\frac{\lambda}{D} = 4.44 \times 10^{-8}\,\mathrm{rad}
     = 9.16\,\mathrm{mas} \]
  \hfill(1.0)

  Now let us derive the time dependence of the angular separation between two
  stars as seen from the earth.

  % FIGURE NEEDED: geometry sketch — two concentric circular orbits about the
  % centre of mass C, stars A and B on a diameter making angle theta with the
  % perpendicular to the line of sight, projected separation s toward the
  % Observer (2022 theory solutions, page 4 of 21).

  As it is seen from the figure, distance $s = l\cos\theta$, where $l$ is the
  distance between stars and $\theta$ is the angle of rotation of the stars
  about their center of mass after the stars are seen with maximum angular
  separation. The angular separation seen by an observer on the earth:
  \[ \gamma = \frac{s}{d} = \frac{l}{d}\cos\theta \]
  \hfill(1.0)

  The 38 days, during which objects are seen as a one object the system rotates
  by the angle $\beta = 2\pi\,\frac{38}{100}$. \hfill(1.0)

  Thus, the corresponding starting angle $\theta'$ at which stars are no longer
  distinguishable:
  \[ \theta' = \frac{\pi}{2} - \frac{\beta}{4}
     = \frac{\pi}{2} - \frac{19\pi}{100} = \frac{31\pi}{100} \]
  \hfill(1.5)

  and this is the moment, when stars stop to be distinguishable. Thus,
  corresponding angular separation on the sky is $\delta$:
  \[ \delta = \frac{l}{d}\cos\theta' \]
  From which one obtains
  \[ l = \frac{\delta d}{\cos\theta'} = 1.45\,\mathrm{au} \]
  \hfill(1.5)

\item[(c)] From the Kepler's 3rd law:
  \begin{align*}
    T^2 &= \frac{4\pi^2 l^3}{G(M_A + M_B)} \\
    M_A + M_B &= \frac{4\pi^2 l^3}{G T^2} = 40.7\,M_\odot
  \end{align*}
  \hfill(2.0)

\item[(d)] From given data for case 1:
  \begin{align*}
    (U - V)_1 &= (U - B)_1 + (B - V)_1 = 0.3 \\
    U_1 &= U_{0_1} + a_U d = 6.51 \\
    V_1 &= U_1 - (U - V)_1 = 6.21 \\
    m_{Bol_1} &= V_1 + BC_1 = 6.31
  \end{align*}
  \hfill(1.5)

  Similarly for second configuration:
  \begin{align*}
    (U - V)_2 &= (U - B)_2 + (B - V)_2 = 0.37 \\
    U_2 &= U_{0_2} + a_U d = 6.98 \\
    V_2 &= U_2 - (U - V)_2 = 6.61 \\
    m_{Bol_2} &= V_2 + BC_2 = 6.79
  \end{align*}
  \hfill(1.0)

  The difference in bolometric magnitudes:
  \begin{align*}
    m_{bol_2} - m_{bol_1}
      &= -2.5\log\left(\frac{L_2}{L_1}\right) \\
      &= -2.5\log\left(
         \frac{\pi(R_A^2 - R_B^2)\sigma T_A^4 + \pi R_B^2 \sigma T_B^4}
              {\pi R_A^2 \sigma T_A^4 + \pi R_B^2 \sigma T_B^4}\right)
  \end{align*}
  \hfill(1.5)
  \begin{align*}
    \therefore\ 6.79 - 6.31
      &= -2.5\log\left(
         \frac{(R_A^2 - R_B^2)T_A^4 + R_B^2 T_B^4}
              {R_A^2 T_A^4 + R_B^2 T_B^4}\right) \\
      &= -2.5\log\left(
         \frac{(R_A^2/R_B^2 - 1)T_A^4 + T_B^4}
              {T_A^4 R_A^2/R_B^2 + T_B^4}\right) \\
    \therefore\ 0.48
      &= -2.5\log\left[
         \frac{\left(\frac{R_A^2}{R_B^2} - 1\right) + \frac{T_B^4}{T_A^4}}
              {\frac{T_B^4}{T_A^4} + \frac{R_A^2}{R_B^2}}\right]
  \end{align*}
  \hfill(1.0)
  \[ \therefore\ 10^{\frac{-0.48}{2.5}}
     = \frac{\frac{R_A^2}{R_B^2} - 1 + \frac{T_B^4}{T_A^4}}
            {\frac{T_B^4}{T_A^4} + \frac{R_A^2}{R_B^2}}
     = 10^{-0.19} = 0.65 \]
  \hfill(1.0)
  \begin{align*}
    \frac{R_A^2}{R_B^2} - 1 + \frac{T_B^4}{T_A^4}
      &= 0.65\,\frac{T_B^4}{T_A^4} + 0.65\,\frac{R_A^2}{R_B^2} \\
    0.35\,\frac{R_A^2}{R_B^2}
      &= 1 - 0.35\,\frac{T_B^4}{T_A^4} \\
      &= 1 - 0.35 \times 0.482 = 0.83 \\
    \therefore\ \frac{R_A}{R_B}
      &= \sqrt{\frac{0.83}{0.35}} = 1.54
  \end{align*}
  \hfill(1.0)
  \[ \frac{M_A}{M_B} = \frac{\rho_A R_A^3}{\rho_B R_B^3}
     = 1.54^3 \cdot 0.7 = 2.56 \]
  \hfill(1.0)
  \[ \text{But, } M_A + M_B = 40.7\,M_\odot \]
  \hfill(1.0)
  \[ \therefore\ M_B = \frac{40.7\,M_\odot}{3.56} = 11.4\,M_\odot \]
  \hfill(1.0)
  \[ M_B = 29.3\,M_\odot \]
  % sic — the source writes M_B again; this last line is clearly M_A
  % (= 40.7 - 11.4 = 29.3 solar masses).
\end{description}
